\documentclass{article}
\usepackage{iclr2026_conference,times}
\input{math_commands.tex}

\usepackage{hyperref}
\usepackage{url}
\usepackage{amsmath}
\usepackage{amssymb}
\usepackage{booktabs}
\usepackage{xcolor}



% -----------------------------------------------------------------------------
% Internal result placeholders. Replace every \Run{...} before submission.
% Do not submit an ICLR paper with these unresolved.
% -----------------------------------------------------------------------------
\newcommand{\Run}[1]{\textcolor{red}{\textbf{[RUN: #1]}}}

\title{BGFM: Boltzmann-Calibrated Flow Matching with\\
Universal Neural Potentials for 3D Molecular Generation}

\author{Anonymous Authors\\
Under double-blind review}

\begin{document}
\maketitle

\begin{abstract}
Fast 3D molecular flow and diffusion models generate diverse de novo
molecules, but their probability geometry is learned from empirical
molecular datasets rather than from a physical energy landscape.
Energy-based molecular generators recover an explicit scalar energy,
but sample with expensive Langevin or parallel-tempering chains and can
trade diversity for physical quality. We propose \textbf{BGFM}, an
external-neural-potential calibrated flow generator for de novo 3D
molecular generation. BGFM keeps a FlowMol3-style amortized proposal,
but trains it with OMol25-derived forces and energies~\citep{omol25} through two
Boltzmann regularizers: late-time force-score consistency and
within-parent density-energy variance. BGFM also learns a calibrated
scalar energy head from the same external neural-potential teacher and
uses it for a short coordinate-only learned-energy corrector after the
flow proposal. The resulting model is designed to combine the speed and
diversity of flow matching with the low-strain sampling behavior of
energy-based models. Because OMol25 labels enter training, we separate
mechanistic OMol25 alignment from external validation: the headline
experiments use standard QM9/GEOM-Drugs generation metrics, independent
GFN2-xTB/MMFF/DFT relaxation oracles, quality-diversity-compute Pareto
curves, and shuffled-label negative controls. This paper should be
submitted only after replacing all marked result placeholders with
completed runs.
\end{abstract}

\section{Introduction}

Three-dimensional molecular generators are increasingly used as the
front end of simulation and discovery pipelines. A useful generated
structure must not only be chemically valid; it should already be close
to a physically plausible region of an energy landscape, otherwise
substantial relaxation or molecular dynamics is required before the
sample can be used. Current models expose a tradeoff. Flow-matching and
diffusion models amortize sampling and produce diverse molecules, but
standard training matches empirical data distributions and does not, by
construction, enforce Boltzmann relative probabilities under a
quantum-level energy model. Energy-based molecular generators learn an
explicit scalar potential and can rank, filter, and steer samples, but
sampling from a learned energy generally requires many inference-time
steps.

BGFM addresses this gap by treating a universal neural potential as an
external physical teacher for an amortized molecular flow generator. The
model keeps a FlowMol3-style de novo proposal~\citep{flowmol3} over atom types, charges,
and coordinates, but adds training-time physical calibration from OMol25.
The first calibration term aligns the flow-implied coordinate score with
neural-potential forces near the data endpoint. The second estimates a
coordinate log-density with a FFJORD integral and enforces that, within a
cloud of perturbations of the same parent molecule,
\begin{equation}
    \log p_\theta^{\rm pos}(r \mid c) + \beta E_{\rm NP}(r,c)
\end{equation}
varies as little as possible. This is a local conditional Boltzmann
condition: lower-energy perturbations of the same molecular identity
should receive higher probability. BGFM further trains a small scalar
energy head to approximate the external neural-potential energy and
force, then uses this head for a short learned-energy Langevin corrector
that refines coordinates while holding the sampled molecular identity
fixed.

The central comparison is with EBMol-style energy-based generation~\citep{ebmol}.
EBMol learns a persistent scalar energy and samples with Mirror-Langevin
dynamics and parallel tempering. BGFM takes the complementary route:
keep the fast flow proposal, calibrate it with external physical labels,
and use a short learned-energy corrector rather than a long MCMC chain.
Therefore, our primary evaluation is not just validity or an OMol25
self-correlation. We evaluate whether BGFM improves the
physical-quality--diversity--compute frontier under independent
relaxation oracles.

\paragraph{Contributions.}
\begin{enumerate}
    \item \textbf{External-oracle calibrated flow generation.} We propose
    a single de novo 3D generator that combines a FlowMol3-style proposal,
    OMol25 force/density Boltzmann regularization, a calibrated energy
    head, and a short learned-energy corrector.
    \item \textbf{Conditional Boltzmann regularization.} BGFM targets the
    coordinate Boltzmann distribution conditional on molecular identity,
    avoiding the ill-defined claim of a global Boltzmann distribution over
    arbitrary atom counts, compositions, and charges.
    \item \textbf{Amortized proposal plus calibrated correction.} Rather
    than replacing the generator with an EBM, BGFM learns a calibrated
    scalar energy head and applies only a small number of refinement steps
    after fast flow sampling.
    \item \textbf{Evaluation protocol against self-confirming metrics.}
    We evaluate standard generative quality, independent physical quality,
    quality-diversity-compute Pareto behavior, mechanism-level Boltzmann
    alignment, ablations, and negative controls.
\end{enumerate}

\section{Related Work}

\paragraph{3D molecular flow and diffusion.}
EDM, GeoLDM, MiDi, GeoBFN, SemlaFlow, FlowMol, and FlowMol3 learn
transport or denoising maps for molecular atom types and coordinates.
FlowMol3~\citep{flowmol3} is the closest backbone baseline: it is an open-source
multi-modal flow-matching model for de novo all-atom small-molecule
generation and improves FlowMol using self-conditioning, fake atoms,
and train-time geometry distortion without changing the underlying GNN
or flow formulation. BGFM is not a replacement backbone; it is an
external-neural-potential calibration layer on top of this family of
amortized generators.

\paragraph{Energy-based molecular generation.}
EBMol~\citep{ebmol} learns a time-unconditional atom-additive scalar potential and
samples it with Mirror-Langevin dynamics plus parallel tempering. It is
therefore the direct competitor for physically motivated 3D generation.
BGFM differs in two ways. First, BGFM uses an external universal neural
potential as the physical teacher, whereas EBMol learns an energy from
geometry data without direct energy supervision. Second, BGFM retains an
amortized flow proposal and uses a short corrector, whereas EBMol relies
on inference-time energy sampling.

\paragraph{Flexible conformer ensemble generation.}
FlexiFlow~\citep{flexiflow} extends flow matching to jointly sample molecular graphs and
multiple conformations and reports strong QM9/GEOM-Drugs performance.
FlexiFlow is a required standard-generation baseline, especially for
conformer-ensemble metrics. Its goal is complementary to BGFM: FlexiFlow
focuses on molecular/conformer ensemble coverage, while BGFM focuses on
external-neural-potential calibrated low-strain de novo generation.

\paragraph{Boltzmann generators and energy-based flow matching.}
Boltzmann Generators, equivariant flow matching for Boltzmann sampling,
and iEFM~\citep{iefm} study sampling from known or learned unnormalized densities,
often for fixed systems or continuous toy/physical potentials. iEFM is
especially relevant theoretically because it trains CNFs from energy
function evaluations, but its evaluation setting is fixed continuous
Boltzmann sampling rather than mixed discrete-continuous de novo molecule
generation.

\paragraph{Universal neural potentials.}
OMol25~\citep{omol25} provides large-scale DFT data and neural-potential models over a
broad chemistry space. BGFM uses such a neural potential not as a
post-hoc relaxer but as a training-time teacher for forces, energies,
and energy-head calibration.

\section{Method}

Let a molecular configuration be $x=(c,r)$, where $c$ denotes molecular
identity variables (atom count, atom types, and charges) and
$r\in\mathbb R^{3N}$ denotes coordinates. The external neural potential
returns an energy $E_{\rm NP}(r,c)$ and forces
$F_{\rm NP}(r,c)=-\nabla_r E_{\rm NP}(r,c)$. BGFM does not define a
global canonical ensemble over arbitrary molecular identities. Instead,
for a fixed identity $c$, it targets the conditional coordinate
Boltzmann law
\begin{equation}
    p_B(r\mid c,T)=\frac{1}{Z_c(T)}\exp[-\beta E_{\rm NP}(r,c)],
    \qquad \beta=1/(kT).
    \label{eq:conditional-boltzmann}
\end{equation}
This target is well-defined because $c$ is fixed; cross-composition
probabilities are left to the empirical proposal distribution and are
not interpreted as thermodynamic weights.

\subsection{Flow proposal}

BGFM starts from a FlowMol3-style flow proposal. Continuous coordinates
use a Gaussian prior and a time-dependent velocity field
$v_\theta(x_t,t)$; atom types and charges use the categorical flow/CTMC
component of the backbone. The base training loss is the standard
flow-matching objective
\begin{equation}
    \mathcal L_{\rm FM}=\mathbb E\left[\|v_\theta(x_t,t)-u_t\|^2 +
    \mathcal L_{\rm disc}\right].
\end{equation}
The proposal is responsible for chemical diversity and fast amortized
sampling. The physical modules below calibrate this proposal rather than
replace it.

\subsection{Late-time force-score consistency}

For the conditional Boltzmann law in Eq.~\eqref{eq:conditional-boltzmann},
\begin{equation}
    \nabla_r \log p_B(r\mid c,T)=\beta F_{\rm NP}(r,c).
\end{equation}
Under the linear Gaussian interpolant, the coordinate score at time $t$
can be estimated from the flow velocity by
\begin{equation}
    s_\theta(x_t,t)=\frac{t\,v_\theta(x_t,t)-x_t}{(1-t)\sigma_0^2}.
\end{equation}
BGFM evaluates this score at late but finite times
$t\in\mathcal T$ and aligns it with the endpoint neural-potential force:
\begin{equation}
    \mathcal L_{\rm force}=\mathbb E_{t\in\mathcal T}
    \left[1-\cos\left(s_\theta(x_t,t),\,\beta F_{\rm NP}(r_1,c)\right)\right].
    \label{eq:force-loss}
\end{equation}
This is a late-time regularizer, not an exact identity at arbitrary
intermediate $x_t$. It avoids online neural-potential calls during every
training step by using precomputed endpoint forces.

\subsection{Density-energy variance}

The most direct probability calibration is that
$\log p_B(r\mid c,T)+\beta E_{\rm NP}(r,c)$ equals a constant depending
only on $c$ and $T$. BGFM estimates the coordinate log-density
$\log p_\theta^{\rm pos}(r\mid c)$ with a FFJORD-style divergence
integral through the coordinate velocity field. For each parent molecule
$m$, we precompute a cloud of perturbations
$\{r_m^{(k)}\}_{k=1}^K$ and neural-potential energies. The loss is
\begin{equation}
    \mathcal L_{\rm dens}=\mathbb E_m\operatorname{Var}_{k}
    \left[
        \log p_\theta^{\rm pos}(r_m^{(k)}\mid c_m)
        +\beta E_{\rm NP}(r_m^{(k)},c_m)
    \right].
    \label{eq:density-energy}
\end{equation}
Because the variance is taken within the same parent molecule, the
unknown offset $-\log Z_{c_m}(T)$ cancels. We optionally add a
composition-conditioned scalar offset $b_\phi(c)$ as an anchor, but do
not interpret it as a supervised thermodynamic partition-function
estimator.

\subsection{Calibrated energy head}

To obtain an explicit inference-time energy without running OMol25 at
sampling time, BGFM trains a scalar energy head $\hat E_\psi(r,c)$ on the
same external labels:
\begin{equation}
    \mathcal L_{\rm head}=\|\hat E_\psi(r,c)-E_{\rm NP}(r,c)\|_1
    +\lambda_F\| -\nabla_r\hat E_\psi(r,c)-F_{\rm NP}(r,c)\|_2^2.
    \label{eq:head-loss}
\end{equation}
The head is a calibrated strain scorer used for ranking and refinement.
It is not claimed to be a free energy.

\subsection{Short learned-energy corrector}

After the flow proposal samples $(c,r^0)$, BGFM optionally refines only
coordinates using the calibrated head:
\begin{equation}
    r^{j+1}=\Pi\left[r^j+\eta_j\left(-\nabla_r\hat E_\psi(r^j,c)\right)
    +\sqrt{2\eta_j/\beta}\,\epsilon_j\right],\qquad j=0,\ldots,J-1.
    \label{eq:corrector}
\end{equation}
Here $J\in\{0,20,50,100\}$ in our compute sweep, $\Pi$ removes the
center-of-mass drift, and discrete identity $c$ is fixed. We report
flow NFEs, energy-head NFEs, external oracle calls, wall-clock time, and
failure rate. Unless a Metropolis correction is explicitly enabled and
reported, this corrector is an unadjusted short refinement, not an exact
MCMC sampler.

\subsection{Training objective}

The full training loss is
\begin{equation}
    \mathcal L = \mathcal L_{\rm FM}
    +\lambda_1\mathcal L_{\rm force}
    +\lambda_2\mathcal L_{\rm dens}
    +\lambda_3\mathcal L_{\rm anchor}
    +\lambda_4\mathcal L_{\rm head}.
\end{equation}
All auxiliary weights use warmup/ramp schedules. OMol25 is not called in
the main training loop: endpoint forces are stored in the dataset and
perturbation energies are precomputed offline.

\paragraph{Optional joint-density extension.}
BGFM can also add discrete log-probability terms for atom counts, atom
types, and charges to form
$\log p_\theta(x)=\log p_\theta^{\rm pos}(r\mid c)+\log p_\theta(c)$.
This extension is useful for diagnostics and future full-joint density
calibration. Unless explicitly enabled and verified in an ablation, the
main objective in this paper should be interpreted as conditional
coordinate Boltzmann calibration.

\section{Experiments}

\paragraph{Experimental questions.}
We design experiments to answer six questions. First, does BGFM preserve
standard de novo 3D generation quality? Second, does it improve physical
quality under independent evaluators rather than only under the OMol25
labels used for training? Third, does it achieve a better
quality-diversity-compute frontier than EBMol-style energy sampling?
Fourth, does the proposed objective actually produce local conditional
Boltzmann alignment? Fifth, do the gains require correct physical
labels rather than arbitrary regularization? Sixth, does the calibrated
energy head rank and filter molecular strain better than existing energy
scorers?

\paragraph{Datasets and baselines.}
We evaluate on QM9, GEOM-Drugs, and held-out OMol25 chemistry slices.
Standard generation baselines include EDM, GeoLDM/GeoBFN/SemlaFlow,
FlowMol3 or the same-backbone FM baseline, FlexiFlow, and EBMol. We
include BGFM ablations: flow-only, force-only, density-only, head-only,
without corrector, and full BGFM.

\subsection{Standard generation quality}

We generate 10k samples per model under matched or reported NFE budgets
and report atom stability, molecule stability, validity, uniqueness,
novelty, Vendi diversity, NFE, and wall-clock time.

\begin{table}[h]
\centering
\caption{Standard de novo 3D generation quality. Replace all placeholders before submission.}
\begin{tabular}{lrrrrrrr}
\toprule
Model & NFE & Atom stab. & Mol. stab. & Valid & Unique & Novel & Vendi \\
\midrule
FlowMol3/FM & \Run{NFE} & \Run{} & \Run{} & \Run{} & \Run{} & \Run{} & \Run{} \\
FlexiFlow & reported & \Run{} & \Run{} & \Run{} & \Run{} & \Run{} & \Run{} \\
EBMol & matched & \Run{} & \Run{} & \Run{} & \Run{} & \Run{} & \Run{} \\
BGFM w/o corrector & matched & \Run{} & \Run{} & \Run{} & \Run{} & \Run{} & \Run{} \\
BGFM full & matched & \Run{} & \Run{} & \Run{} & \Run{} & \Run{} & \Run{} \\
\bottomrule
\end{tabular}
\end{table}

\subsection{Independent physical quality}

Because BGFM uses OMol25 labels during training, OMol25 correlation is
not the headline metric. Our main physical-quality table uses independent
GFN2-xTB relaxation, MMFF relaxation on organic subsets, and a DFT
validation subset. We report both valid-connected-only and all-samples
protocols.

\begin{table}[h]
\centering
\caption{Independent physical quality under evaluators not used as BGFM's density-energy target.}
\begin{tabular}{lrrrrrr}
\toprule
Model & Valid conn. & xTB med. $\Delta E$ & xTB mean $\Delta E$ & p90 $\Delta E$ & Relax steps & Failure \\
\midrule
EBMol-low & \Run{} & \Run{} & \Run{} & \Run{} & \Run{} & \Run{} \\
EBMol-high & \Run{} & \Run{} & \Run{} & \Run{} & \Run{} & \Run{} \\
FlowMol3/FM & \Run{} & \Run{} & \Run{} & \Run{} & \Run{} & \Run{} \\
BGFM w/o head & \Run{} & \Run{} & \Run{} & \Run{} & \Run{} & \Run{} \\
BGFM w/o corrector & \Run{} & \Run{} & \Run{} & \Run{} & \Run{} & \Run{} \\
BGFM full & \Run{} & \Run{} & \Run{} & \Run{} & \Run{} & \Run{} \\
\bottomrule
\end{tabular}
\end{table}

\subsection{Quality-diversity-compute Pareto}

We sweep flow NFEs and corrector steps $J\in\{0,20,50,100\}$ and compare
to EBMol's compute-scaled samplers. The key metric is the number of
valid, unique, low-strain samples per GPU-hour:
\begin{equation}
    \mathrm{VULS}_\tau=\frac{\#\{\text{valid, unique samples with }\Delta E_{\rm xTB}<\tau\}}{\text{GPU-hour}}.
\end{equation}
This evaluation targets the main BGFM hypothesis: a fast amortized flow
proposal plus a short calibrated corrector should dominate long
energy-sampling chains on physical quality, diversity, and cost.

\subsection{Mechanistic Boltzmann alignment}

For held-out parent molecules and perturbation families, we measure the
within-parent correlation and calibration slope between
$\log p_\theta^{\rm pos}(r\mid c)$ and $-\beta E(r,c)$. We report this
with OMol25 as a mechanism diagnostic, and repeat it with xTB and a DFT
subset to test transfer.

\subsection{Ablations and negative controls}

A valid physical improvement should disappear or weaken under corrupted
labels. We therefore train or evaluate controls with shuffled forces,
shuffled energies, wrong parent grouping, wrong temperature, force norm
only, and random scalar regularization.

\begin{table}[h]
\centering
\caption{Negative controls separating correct physical supervision from arbitrary regularization.}
\begin{tabular}{lrrrrr}
\toprule
Model/control & OMol25 align. & xTB med. $\Delta E$ & DFT subset & Valid & Vendi \\
\midrule
FM baseline & \Run{} & \Run{} & \Run{} & \Run{} & \Run{} \\
BGFM true labels & \Run{} & \Run{} & \Run{} & \Run{} & \Run{} \\
Shuffled force & \Run{} & \Run{} & \Run{} & \Run{} & \Run{} \\
Shuffled energy & \Run{} & \Run{} & \Run{} & \Run{} & \Run{} \\
Wrong parent grouping & \Run{} & \Run{} & \Run{} & \Run{} & \Run{} \\
Wrong $kT$ & \Run{} & \Run{} & \Run{} & \Run{} & \Run{} \\
\bottomrule
\end{tabular}
\end{table}

\subsection{Cross-model ranking and filtering}

We collect samples from EDM/FlowMol3/EBMol/BGFM and deliberately strained
perturbations. We compare EBMol energy, OMol25 energy, BGFM's calibrated
energy head, and the BGFM combined score as predictors of independent
xTB/DFT relaxation energy and high-strain failure. Metrics are Pearson,
Spearman, AUROC for high-strain detection, and top-decile filtering
quality.

\section{Limitations}

BGFM inherits OMol25's coverage and accuracy limitations. The main
density loss is conditional-coordinate unless the optional discrete
joint-density extension is explicitly enabled. The short corrector is a
learned-energy refinement, not an exact MCMC sampler. The energy head is
a calibrated scorer distilled from neural-potential energies and forces,
not a free-energy estimator. These limitations are evaluated with
independent oracles and shuffled-label controls rather than hidden.

\section{Conclusion}

BGFM is designed to bridge two strengths: amortized de novo flow
sampling and physically meaningful energy landscapes. By calibrating a
fast molecular flow proposal with an external universal neural potential
and using a short learned-energy corrector, BGFM targets low-strain,
diverse, computationally efficient 3D molecule generation. The decisive
empirical question is not whether BGFM improves an OMol25 self-metric,
but whether it improves independent physical quality and the
quality-diversity-compute frontier against strong flow and energy-based
baselines.

\begin{thebibliography}{9}

\bibitem[Griesbacher et~al.(2026)]{ebmol}
Christoph Griesbacher, Lea Bogensperger, Andreas Habring, and Thomas Pock.
\newblock Generating physically consistent molecules with energy-based models.
\newblock \emph{arXiv:2605.18381}, 2026.

\bibitem[Tedoldi et~al.(2025)]{flexiflow}
Riccardo Tedoldi, Ola Engkvist, Patrick Bryant, Hossein Azizpour, Jon Paul Janet, and Alessandro Tibo.
\newblock FlexiFlow: decomposable flow matching for generation of flexible molecular ensemble.
\newblock \emph{arXiv:2511.17249}, 2025.

\bibitem[Woo and Ahn(2024)]{iefm}
Dongyeop Woo and Sungsoo Ahn.
\newblock Iterated energy-based flow matching for sampling from Boltzmann densities.
\newblock \emph{arXiv:2408.16249}, 2024.

\bibitem[Dunn and Koes(2025)]{flowmol3}
Ian Dunn and David R. Koes.
\newblock FlowMol3: Flow matching for 3D de novo small-molecule generation.
\newblock \emph{arXiv:2508.12629}, 2025.

\bibitem[Levine et~al.(2025)]{omol25}
Daniel S. Levine, Muhammed Shuaibi, Evan Walter Clark Spotte-Smith, Michael G. Taylor, Muhammad R. Hasyim, Kyle Michel, Ilyes Batatia, G\'abor Cs\'anyi, Misko Dzamba, Peter Eastman, and others.
\newblock The Open Molecules 2025 (OMol25) dataset, evaluations, and models.
\newblock \emph{arXiv:2505.08762}, 2025.

\end{thebibliography}

\end{document}
